\section{Counting on one physical interval}
\label{sec:counting-tools}

The upper bound follows from two elementary counting criteria on a
connected interval. One uses the orientation of displacement critical
points; the other uses the orientation of curvature zeros. Each criterion
counts the whole physical interval, including the regions between the
captured roots.

\subsection{A common coordinate and the counting criteria}

Fix the incoming raw height coordinate in Lemma~\ref{lem:closing} before
normalizing the tails. Absorb the regular maps between the fixed raw ports
and the chosen sections into the actual adapters. All regional estimates
will then concern the displacement
\begin{equation}\label{sharp:upper-displacement}
 F(x)=x-m\Pi(x),\qquad m>0.
\end{equation}
Every local estimate below includes the divided adapter
\eqref{eq:divided-closing}. Its derivatives remain bounded, with positive
first derivative, as \(m\downarrow0\). Apply the local tail estimates
with these adapters and normalize each map by its exact center and first
variation. The relation with the physical map is the identity
\begin{equation}\label{sharp:upper-affine}
 \Pi(a_++c_+w)=a_-+c_-\widehat\Pi(w),\qquad c_\pm>0.
\end{equation}
Thus the signs of second and third derivatives transfer by the positive
factors \(c_-/c_+^2\) and \(c_-/c_+^3\). The affine relation is
essential here: independently chosen nonlinear coordinates need not
preserve curvature. Comparisons obtained in an inverse coordinate are
therefore transferred to the direct displacement before counting its
critical points.

All carriers below are intervals of actual through trajectories in the
nested section windows. Scalar first variation orders them by the central
height. The preimage of an output interval under an increasing passage map
is an interval. Consequently the interval between any finite set of
captured roots stays in the same enlarged physical domain. The estimates
used between roots therefore apply at nonroot points as well.

\begin{lemma}\label{sharp:upper-critical-count}
Let \(F\) be \(C^3\) on a connected carrier \(I\). If every critical
point satisfies \(F''>0\), then \(F\) has at most one critical point,
which is simple and upward, and at most two isolated zeros, counted with
multiplicity. More generally, let \(J\subseteq I\) be a connected interval,
possibly empty or all of \(I\). If \(F'\) has at most \(j\)
zeros on \(J\), counted with multiplicity, where \(0\le j\le2\), and \(F'>0\) on
\(I\setminus J\), then \(F'\) has at most \(j\) zeros on \(I\), and
\(F\) has at most \(j+1\) isolated zeros there, counted with multiplicity.
For \(J=\varnothing\), take \(j=0\): there is no critical point and
at most one simple displacement zero.
\end{lemma}
\begin{proof}
In the first case every zero of \(F'\) is a simple upward crossing.
Two such crossings would require an intervening zero with another
orientation, so there is at most one. In the second case no critical
point lies outside \(J\), whether its complement is empty or lies on
one or both sides. Rolle's theorem with
multiplicities proves both assertions.
\end{proof}

\begin{lemma}\label{sharp:upper-curvature-count}
Let \(F=x-m\Pi\), \(m>0\), be \(C^3\) on a connected carrier \(I\),
and let \(J\subseteq I\) be a connected interval, possibly empty or
all of \(I\). Suppose that every zero of \(\Pi''\) on \(J\) is simple
and has the same crossing orientation, and that \(F'>0\) on
\(I\setminus J\). Then \(\Pi''\) has at most one zero on \(J\),
\(F'\) has at most two zeros on \(I\), counted with multiplicity,
and \(F\) has at most three
isolated displacement zeros on \(I\), counted with multiplicity.
If \(J=\varnothing\), there is no critical point and at most one
simple displacement zero.
\end{lemma}
\begin{proof}
For nonempty \(J\), two simple curvature zeros of the same orientation
would require an intervening zero with another orientation, so there is
at most one. Since \(F''=-m\Pi''\), Rolle's theorem with multiplicities
gives at most two zeros of \(F'\) on \(J\). Apply
Lemma~\ref{sharp:upper-critical-count} on \(I\). For empty \(J\),
\(F'>0\) throughout \(I\). Positive affine changes in
\eqref{sharp:upper-affine} preserve the curvature orientation, even if
their scale ratios degenerate.
\end{proof}

The regional constructions specify the actual intervals \(I,J\) and
establish the hypotheses at every point of those intervals. For each
coefficient fibre, apply Lemma~\ref{sharp:upper-critical-count} or
Lemma~\ref{sharp:upper-curvature-count} to its whole carrier. The later
coefficient cover takes the maximum of these bounds; its overlap choices
are independent of the phase input.

\subsection{Ordinary adapters and divided closing maps}

The next two lemmas control the regular maps appended to a singular
passage. They use ordinary third derivatives of those maps and allow the
singular passage's first gain to tend to zero.

\begin{lemma}\label{phys:lem-c3-adapters}
Let the ordinary adapters below have uniform \(C^3\) norms and positive
first derivatives bounded away from zero on their original buffered
images. All parameters and gains are frozen under the phase derivation
\(D\). The following identities and bounds are uniform in those
parameters.

For an actual raw output \(T\) in such an image, put
\(\ell=\log\psi'\). Then
\begin{equation}\label{phys:eq-observer-log}
 \begin{aligned}
 D\log\frac{(\psi\circ T)'}{T'}&=\ell'(T)DT,\\
 D^2\log\frac{(\psi\circ T)'}{T'}
 &=\ell''(T)(DT)^2+\ell'(T)D^2T .
 \end{aligned}
\end{equation}
Here \(T'>0\) is the exact raw first gain; the formulas impose no lower
bound on it. If the full translated output and its two absolute phase
derivatives are bounded, and \(mT\) stays in the fixed closing buffer,
then
\begin{equation}\label{phys:eq-absolute-closing-two}
 \sum_{j=0}^2\left(
 |D^j(f(mT)-f(0))|
 +|D^j(\log f'(mT)-\log f'(0))|\right)\le C m.
\end{equation}
The constant depends on these absolute output bounds, not on a
normalized raw gain or a ratio of gains.

For \(f(0)=0\), the divided closing map \(C_m(q)=f(md(q))/m\),
extended at zero by \(C_0=f'(0)d\), has jets
\begin{equation}\label{phys:eq-closing-jets}
 \begin{aligned}
 C_m'&=f'(md)d',\\
 C_m''&=m f''(md)(d')^2+f'(md)d'',\\
 C_m'''&=m^2f'''(md)(d')^3
       +3m f''(md)d'd''+f'(md)d''' .
 \end{aligned}
\end{equation}
They are uniformly bounded and the first is a positive unit on each
prescribed compact \(q\)-interval with buffered \(md\)-image.
Finally, for a coefficient-only center \(z_0\), \(0<d\le1\), and a
fixed compact phase interval whose \(z_0+dy\)-image stays buffered,
\begin{equation}\label{phys:eq-observer-normalization}
 N_d(y)=\frac{\psi(z_0+dy)-\psi(z_0)}{d\psi'(z_0)},\qquad
 N_0(y)=y,\qquad \|N_d-\mathrm{id}\|_{C^3}\le C d.
\end{equation}
This last bound concerns phase derivatives, not derivatives of the
center or scale with respect to a parameter.
\end{lemma}
\begin{proof}
The chain rule proves \eqref{phys:eq-observer-log} and
\eqref{phys:eq-closing-jets}; the integral formula
\(C_m=d\int_0^1f'(tmd)\,dt\) gives the extension at zero.
The ordinary \(C^3\) bounds and positive slopes then give
\eqref{phys:eq-absolute-closing-two} by the mean value theorem.
Taylor's integral formula gives
\eqref{phys:eq-observer-normalization}, whose second and third
rows contain factors \(d\) and \(d^2\), respectively.
\end{proof}

\begin{lemma}\label{phys:lem-c3-core-transport}
Let \(M\) be an intrinsic core map with uniformly bounded \(C^4\)
norm on a fixed buffered phase interval. Let \(A,B\) be ordinary
\(C^3\) adapters close to the identity, with uniformly positive
slopes, on buffers containing all compositions below. On a common
compact interval for \(A^{-1}\) and \(M\circ A^{-1}\),
\begin{equation}\label{phys:eq-c3-core-transport}
 \|B\circ M\circ A^{-1}-M\|_{C^3}
 \le C\bigl(\|A-\mathrm{id}\|_{C^3}
              +\|B-\mathrm{id}\|_{C^3}\bigr).
\end{equation}
The fourth derivative is required only of the intrinsic core.
\end{lemma}
\begin{proof}
The inverse formulas through order three and the uniform slope bound
give \(\|A^{-1}-\mathrm{id}\|_{C^3}\le C\|A-\mathrm{id}\|_{C^3}\).
In the third chain-rule row for \(M\circ A^{-1}-M\), the translated
factor \(M'''(A^{-1})-M'''\) is bounded by
\(\|M''''\|_\infty\|A^{-1}-\mathrm{id}\|_\infty\); the other
rows need fewer core derivatives. The ordinary composition of
\(B-\mathrm{id}\) with \(M\circ A^{-1}\) is bounded in \(C^3\)
without translating any fourth derivative of \(B\). These estimates
prove \eqref{phys:eq-c3-core-transport}.
\end{proof}

\subsection{Weighted curvature of a completed passage}

To use an action calculation in the physical input, we must include the
curvature contributed by each section map. The following identity keeps
these contributions separate and allows a positive weight adapted to
the coefficient chart.

\begin{lemma}[Weighted curvature of a completed passage]
\label{neg:lem-completed-curvature}
Let \(u,v>0\) be radial coordinates, with \(P_\pm,B_\pm>0\)
on the two strict through intervals, and suppose
\[
 v_u=\frac{uB_+(u)P_-(v)}{vB_-(v)P_+(u)}(1+\varrho)>0,
 \quad S(-u^2/2)=-v^2/2,\quad
 \Pi=\phi_-\circ S\circ\phi_+^{-1},
\]
Here \(P_\pm,B_\pm\) are \(C^2\), \(\varrho=\varrho(u,v)\)
is \(C^2\) on the product of the two radial intervals, and
\(\phi_\pm\) are the separately normalized actual \(C^3\) section
maps with positive derivatives. Put
\(\mathcal D_+=P_+\partial_u/(uB_+)\),
\(\mathcal D_-=P_-\partial_v/(vB_-)\), and \(R_B=B_+/B_-\).
For any smooth positive weight \(W=W(u,v)\), one has exactly
\begin{align}
 \phi_+'\frac{P_+}{W}\frac{\Pi''}{\Pi'}
 ={}&\frac{P_+' /u-R_B(1+\varrho)P_-'/v}{W}
 -\frac{P_+}{uW}\frac d{du}\log R_B \notag\\
 &-\frac{B_+}{W}
       \{\mathcal D_++(1+\varrho)\mathcal D_-\}\log(1+\varrho)
       \notag\\
 &-\frac{B_+(1+\varrho)}{W}\mathcal D_-\log\phi_-'
 +\frac{P_+}{u^2W}(u\partial_u)\log\phi_+'.
 \label{deg:weighted-exact-curvature-chain}
\end{align}
Primes on each \(\phi_\pm\) use its own quadratic label; primes on
\(P_\pm\) use the displayed radial variable. The derivative \(d/du\)
follows the matched graph. In particular both section terms remain in
the curvature of the actual affine input.
\end{lemma}
\begin{proof}
One has \(S'=R_BP_-(1+\varrho)/P_+\). Logarithmically differentiate
this identity using \(d/db_+=-u^{-1}d/du\), and substitute in
\[
 \phi_+'\frac{\Pi''}{\Pi'}=
 \frac{S''}{S'}+S'\frac{\phi_-''}{\phi_-'}-
                         \frac{\phi_+''}{\phi_+'}.
\]
Multiplication by \(P_+/W\) gives the formula. Its coefficients need
not extend without weights to a stopped endpoint.
\end{proof}


